\solutionchapter{Logique}

%────────────────────────────────────────────────────────────────

\subsection*{Exercices}
\addcontentsline{toc}{subsection}{Exercices}

\solutiontitle{Propositions ou non}

\begin{enumerate}
  \item \textbf{Proposition vraie.}\enspace $5\times4=20$.
  \item \textbf{Proposition fausse.}\enspace $\sqrt{2}\approx1{,}414\notin\N$.
  \item \textbf{Non~: pr\'edicat.}\enspace La valeur de $x+3=7$ d\'epend de $x$.
  \item \textbf{Proposition vraie.}\enspace Tout triangle \'equilat\'eral a deux c\^ot\'es \'egaux.
  \item \textbf{Proposition vraie.}\enspace $0\times n=0$ pour tout $n$.
  \item \textbf{Non~: question.}\enspace Une question n'a pas de valeur de v\'erit\'e.
  \item \textbf{Proposition vraie.}\enspace $\pi\approx3{,}14159>3$.
  \item \textbf{Non~: pr\'edicat.}\enspace La valeur d\'epend de $n$.
  \item \textbf{Proposition vraie.}\enspace $\N\subset[0,+\infty)$ par d\'efinition.
  \item \textbf{Non~: ordre.}\enspace Un imp\'eratif n'est ni vrai ni faux.
\end{enumerate}

%────────────────────────────────────────────────────────────────
\solutiontitle{N\'egations simples}

\begin{enumerate}
  \item $17$ n'est \textbf{pas} un nombre premier. (Fausse~: $17$ est premier.)
  \item Le triangle $ABC$ n'est \textbf{pas} rectangle.
  \item $x < 5$.
  \item Il ne pleut pas \textbf{ou} il ne fait pas froid.
        (De Morgan~: $\lnot(P\land Q)\equiv\lnot P\lor\lnot Q$.)
  \item Il ne pleut pas \textbf{et} il ne fait pas chaud.
        ($\lnot(P\lor Q)\equiv\lnot P\land\lnot Q$.)
  \item $\exists n\in\N,\ n^2<n$. (Fausse~: $n(n-1)\geq0$ pour tout $n\in\N$.)
  \item $\forall x\in\R,\ x^2\ne-1$. (Vraie.)
  \item $\exists n\in\Z,\ 2n$ est impair. (Fausse.)
\end{enumerate}

%────────────────────────────────────────────────────────────────
\solutiontitle{Tables de v\'erit\'e}

\renewcommand{\arraystretch}{1.25}
\begin{center}
\begin{tabular}{|cc|cccccc|}
\hline
\rowcolor{BN!12}
$P$ & $Q$ & $\lnot P\lor Q$ & $P\land\lnot Q$ & $(P\lor Q)\land\lnot P$ &
$\lnot(P\land Q)$ & $\lnot P\land\lnot Q$ & $\lnot P\lor\lnot Q$ \\
\hline
\vrai & \vrai & \vrai & \faux & \faux & \faux & \faux & \faux \\
\vrai & \faux & \faux & \vrai & \faux & \vrai & \faux & \vrai \\
\faux & \vrai & \vrai & \faux & \vrai & \vrai & \faux & \vrai \\
\faux & \faux & \vrai & \faux & \faux & \vrai & \vrai & \vrai \\
\hline
\end{tabular}
\end{center}
Remarquer~: la formule~(4) $\lnot(P\land Q)$ est identique \`a~(6) $\lnot P\lor\lnot Q$~:
c'est la premi\`ere loi de De~Morgan.
De m\^eme,~(5) $\lnot P\land\lnot Q$ est la n\'egation de $P\lor Q$.

%────────────────────────────────────────────────────────────────
\solutiontitle{Valeur de v\'erit\'e d'une implication}

\begin{enumerate}
  \item \textbf{Vraie.}\enspace $P=\vrai$, $Q=\vrai$~: $P\Rightarrow Q$ est vraie.
  \item \textbf{Vraie.}\enspace $P=\faux$~: une implication \`a hypoth\`ese fausse est toujours vraie,
        quelle que soit la conclusion.
  \item \textbf{Fausse.}\enspace $P=\vrai$ et $Q=\faux$~: c'est l'unique cas o\`u une implication
        est fausse.
  \item \textbf{Vraie.}\enspace $P$~: \og$n$ divisible par~$4$\fg, $Q$~: \og$n$ divisible par~$2$\fg.
        Si $4\mid n$, alors $n=4k$, donc $n=2(2k)$~: $2\mid n$.
  \item \textbf{Fausse.}\enspace $n=2$ est divisible par $2$, mais pas par $4$.
\end{enumerate}

%────────────────────────────────────────────────────────────────
\solutiontitle{Quantificateurs~: traduction}

\begin{enumerate}
  \item \og Pour tout entier naturel $x$, $x+1>x$.\fg\enspace \textbf{Vraie.}
  \item \og Il existe un entier relatif dont le carr\'e vaut $-4$.\fg\enspace
        \textbf{Fausse}~: un carr\'e est toujours $\geq0$.
  \item \og Pour tout r\'eel $x$, $x^2\geq0$.\fg\enspace \textbf{Vraie.}
  \item \og Il existe un rationnel dont le carr\'e vaut $2$.\fg\enspace
        \textbf{Fausse.} ($\sqrt{2}\notin\Q$.)
  \item \og Pour tout r\'eel non nul, $1/x>0$.\fg\enspace
        \textbf{Fausse.} Contre-exemple~: $x=-1$.
  \item \og Il existe $n\in\N$ tel que $n^2=2n$.\fg\enspace
        \textbf{Vraie.} $n=0$ ou $n=2$ conviennent.
\end{enumerate}

%────────────────────────────────────────────────────────────────
\solutiontitle{\'Ecrire avec des quantificateurs}

\begin{enumerate}
  \item Pour tout quadrilat\`ere $Q$, si $Q$ est un carr\'e, alors $Q$ est un rectangle.
        En abr\'eg\'e~:
        \[\forall Q,\quad Q\text{ carr\'e}\Rightarrow Q\text{ rectangle}.\]
  \item $\exists n\in\N,\ (n \text{ est pair et } n \text{ est premier})$.
        (Vrai~: $n=2$.)
  \item $\forall a,b\in\Z,\ (2\mid a\ \text{et}\ 2\mid b)\Rightarrow 2\mid(a+b)$.
  \item $\forall x\in\R,\ \exists y\in\R,\ x+y=0$. (On prend $y=-x$.)
  \item $\lnot\bigl(\exists n\in\Z,\ 0<n<1\bigr)$, autrement dit
        $\forall n\in\Z,\ n\leq0\ \text{ou}\ n\geq1$.
\end{enumerate}

%────────────────────────────────────────────────────────────────
\solutiontitle{N\'egation avec quantificateurs}

\begin{enumerate}
  \item $\exists x\in\R,\ x^2\leq0$. \textbf{Vraie}~: $x=0$ donne $x^2=0$.
        (La proposition originale est donc fausse.)
  \item $\forall n\in\N,\ 2n+1\ne0$. \textbf{Vraie.}
        (L'originale est fausse~: $2n+1\geq1>0$.)
  \item $\exists n\in\N,\ n^2<n$. \textbf{Fausse.}
        (L'originale est vraie~: $n^2-n=n(n-1)\geq0$.)
  \item $\forall x\in\R,\ x^3\ne-8$. \textbf{Fausse.}
        (L'originale est vraie~: $x=-2$ convient.)
  \item $\exists x\in\R,\ \exists y\in\R,\ x+y\ne y+x$. \textbf{Fausse.}
        (L'originale est vraie~: l'addition est commutative dans $\R$.
        La n\'egation est donc fausse.)
\end{enumerate}

%────────────────────────────────────────────────────────────────
\solutiontitle{R\'eciproque et contrapos\'ee}

\begin{enumerate}
  \item Implication~: si $ABC$ est \'equilat\'eral, alors il est isoc\`ele.
        \textbf{Vraie.} Contrapos\'ee~: si $ABC$ n'est pas isoc\`ele, il n'est pas \'equilat\'eral.
        R\'eciproque~: si $ABC$ est isoc\`ele, il est \'equilat\'eral. \textbf{Fausse}
        (un triangle 5-5-4 est isoc\`ele non \'equilat\'eral).
  \item Implication~: $x=3\Rightarrow x^2=9$. \textbf{Vraie.}
        Contrapos\'ee~: $x^2\ne9\Rightarrow x\ne3$. Vraie.
        R\'eciproque~: $x^2=9\Rightarrow x=3$. \textbf{Fausse}~: $x=-3$ convient.
  \item Implication~: $6\mid n\Rightarrow3\mid n$. \textbf{Vraie.}
        Contrapos\'ee~: $3\nmid n\Rightarrow6\nmid n$. Vraie.
        R\'eciproque~: $3\mid n\Rightarrow6\mid n$. \textbf{Fausse}~: $n=3$ convient.
  \item Implication~: $x>1\Rightarrow x^2>1$. \textbf{Vraie.}
        Contrapos\'ee~: $x^2\leq1\Rightarrow x\leq1$. Vraie.
        R\'eciproque~: $x^2>1\Rightarrow x>1$. \textbf{Fausse}~: $x=-2$ convient.
\end{enumerate}

%────────────────────────────────────────────────────────────────
\solutiontitle{Contre-exemples}

\begin{enumerate}
  \item $n=0$~: $0^2=0\not>0$. (Ou $n=1$~: $1=1\not>1$.)
  \item $x=-3$~: $\sqrt{(-3)^2}=\sqrt{9}=3\ne-3$.
  \item $x=-2$~: $(-2)^2=4$ mais $x\ne2$.
  \item $n=4$~: $4$ est un multiple de $4$, mais $4/8=0{,}5\notin\Z$.
  \item $n=-1$~: $|-1|=1\ne-1$.
\end{enumerate}

%────────────────────────────────────────────────────────────────
\solutiontitle{Implication ou \'equivalence~?}

\begin{enumerate}
  \item $x=2\Rightarrow x^2=4$~: \textbf{vraie.}
        $x^2=4\Rightarrow x=2$~: \textbf{fausse} (contre-exemple $x=-2$).
        Conclusion~: $x=2\Rightarrow x^2=4$ mais non $\Leftrightarrow$.
  \item $n$ pair $\Rightarrow n^2$ pair~: vraie ($n=2k\Rightarrow n^2=4k^2$).
        $n^2$ pair $\Rightarrow n$ pair~: vraie (par contrapos\'ee~: $n$ impair $\Rightarrow n^2$ impair).
        Conclusion~: \textbf{\'equivalence} $n$ pair $\Leftrightarrow n^2$ pair.
  \item $6\mid n\Rightarrow 2\mid n$~: vraie.
        $2\mid n\Rightarrow 6\mid n$~: fausse (contre-exemple~: $n=2$, $6\nmid2$).
        Conclusion~: $6\mid n\Rightarrow 2\mid n$ seulement.
  \item $x=3\Rightarrow|x|=3$~: vraie.
        $|x|=3\Rightarrow x=3$~: fausse ($x=-3$ convient).
        Conclusion~: $x=3\Rightarrow|x|=3$ seulement.
  \item $x=0\Rightarrow xy=0$~: vraie.
        $xy=0\Rightarrow x=0$~: fausse (prendre $x=1$, $y=0$).
        Conclusion~: $x=0\Rightarrow xy=0$ seulement.
\end{enumerate}

%────────────────────────────────────────────────────────────────
\solutiontitle{Table~: contrapos\'ee}

\renewcommand{\arraystretch}{1.35}
\begin{center}
\begin{tabular}{|cc|cc|cc|}
\hline
\rowcolor{BN!12}
$P$ & $Q$ & $\lnot Q$ & $\lnot P$
    & $P\Rightarrow Q$ & $\lnot Q\Rightarrow\lnot P$ \\
\hline
\vrai & \vrai & \faux & \faux & \vrai & \vrai \\
\vrai & \faux & \vrai & \faux & \faux & \faux \\
\faux & \vrai & \faux & \vrai & \vrai & \vrai \\
\faux & \faux & \vrai & \vrai & \vrai & \vrai \\
\hline
\end{tabular}
\end{center}
Les colonnes 5 et 6 sont identiques~: $P\Rightarrow Q\equiv\lnot Q\Rightarrow\lnot P$.\hfill$\square$

%────────────────────────────────────────────────────────────────
\solutiontitle{Table~: connecteur $\Leftrightarrow$}

\begin{center}
\renewcommand{\arraystretch}{1.35}
\begin{tabular}{|cc|cc|c|c|}
\hline
\rowcolor{BN!12}
$P$ & $Q$ & $P\Rightarrow Q$ & $Q\Rightarrow P$
    & $(P\Rightarrow Q)\land(Q\Rightarrow P)$ & $P\Leftrightarrow Q$ \\
\hline
\vrai & \vrai & \vrai & \vrai & \vrai & \vrai \\
\vrai & \faux & \faux & \vrai & \faux & \faux \\
\faux & \vrai & \vrai & \faux & \faux & \faux \\
\faux & \faux & \vrai & \vrai & \vrai & \vrai \\
\hline
\end{tabular}
\end{center}
Les colonnes 5 et 6 co\"incident~: $P\Leftrightarrow Q\equiv(P\Rightarrow Q)\land(Q\Rightarrow P)$.\hfill$\square$

%────────────────────────────────────────────────────────────────
\solutiontitle{Calcul avec valeurs donn\'ees}

$P=\vrai$, $Q=\faux$, $R=\vrai$.
\begin{multicols}{2}
\begin{enumerate}
  \item $P\land Q = \vrai\land\faux = \faux$
  \item $P\lor Q = \vrai\lor\faux = \vrai$
  \item $P\Rightarrow Q = \vrai\Rightarrow\faux = \faux$
  \item $Q\Rightarrow P = \faux\Rightarrow\vrai = \vrai$
  \item $(P\lor Q)\Rightarrow R = \vrai\Rightarrow\vrai = \vrai$
  \item $\lnot P\lor(Q\land R)=\faux\lor(\faux\land\vrai)=\faux\lor\faux=\faux$
  \item $P\Leftrightarrow R = \vrai\Leftrightarrow\vrai = \vrai$
  \item $(P\land\lnot Q)\Rightarrow R = (\vrai\land\vrai)\Rightarrow\vrai = \vrai\Rightarrow\vrai = \vrai$
\end{enumerate}
\end{multicols}

%────────────────────────────────────────────────────────────────
\solutiontitle{Forme logique de r\'esultats math\'ematiques}

\begin{enumerate}
  \item $\forall n\in\Z,\ (n \text{ impair})\Rightarrow(n^2 \text{ impair})$.
  \item $\forall$ triangle $ABC$~:
        $(\text{rectangle en }C)\Leftrightarrow(AB^2=AC^2+BC^2)$.
  \item $\forall n\in\Z,\ \bigl((2\mid n)\land(3\mid n)\bigr)\Rightarrow6\mid n$.
  \item $\forall k\in\Z,\ 2\mid(4k+2)$.
\end{enumerate}

%────────────────────────────────────────────────────────────────
\solutiontitle{Priorit\'e des connecteurs}

Les priorit\'es~: $\lnot$ puis $\land$ puis $\lor$ puis $\Rightarrow$ puis $\Leftrightarrow$.

\begin{enumerate}
  \item $\lnot P\land Q\lor R=((\lnot P)\land Q)\lor R$.
        Elle est vraie ssi $R$ est vraie, ou si $P$ est fausse et $Q$ vraie.
  \item $P\land Q\Rightarrow R=(P\land Q)\Rightarrow R$.
        Elle est fausse seulement quand $P$ et $Q$ sont vraies et $R$ fausse.
  \item $P\lor Q\land\lnot R=P\lor\bigl(Q\land(\lnot R)\bigr)$.
        Elle est fausse seulement quand $P$ est fausse et que $Q\land\lnot R$ est fausse.
\end{enumerate}

\renewcommand{\arraystretch}{1.2}
\begin{center}
\small
\begin{tabular}{|ccc|c|c|c|}
\hline
\rowcolor{BN!12}
$P$&$Q$&$R$ & $((\lnot P)\land Q)\lor R$ & $(P\land Q)\Rightarrow R$ & $P\lor(Q\land\lnot R)$ \\
\hline
\vrai&\vrai&\vrai & \vrai & \vrai & \vrai \\
\vrai&\vrai&\faux & \faux & \faux & \vrai \\
\vrai&\faux&\vrai & \vrai & \vrai & \vrai \\
\vrai&\faux&\faux & \faux & \vrai & \vrai \\
\faux&\vrai&\vrai & \vrai & \vrai & \faux \\
\faux&\vrai&\faux & \vrai & \vrai & \vrai \\
\faux&\faux&\vrai & \vrai & \vrai & \faux \\
\faux&\faux&\faux & \faux & \vrai & \faux \\
\hline
\end{tabular}
\end{center}

%────────────────────────────────────────────────────────────────
\solutiontitle{De Morgan~: simplifications}

\begin{enumerate}
  \item $\lnot(\lnot P\land Q)\equiv\lnot(\lnot P)\lor\lnot Q\equiv P\lor\lnot Q$.
  \item $\lnot(P\lor\lnot Q)\equiv\lnot P\land Q$.
  \item $\lnot\bigl((P\Rightarrow Q)\land P\bigr)
        \equiv\lnot(P\Rightarrow Q)\lor\lnot P
        \equiv(P\land\lnot Q)\lor\lnot P$.
        On simplifie par cas~: si $P$ est fausse, les deux formules
        $(P\land\lnot Q)\lor\lnot P$ et $\lnot P\lor\lnot Q$ sont vraies~;
        si $P$ est vraie, elles valent toutes les deux $\lnot Q$. Donc
        $\lnot\bigl((P\Rightarrow Q)\land P\bigr)\equiv\lnot P\lor\lnot Q$.
        (Autre lecture~: $(P\Rightarrow Q)\land P\equiv P\land Q$, dont la n\'egation est
        $\lnot P\lor\lnot Q$ par De~Morgan.)
  \item $\lnot(\lnot P\lor\lnot Q)\equiv P\land Q$.
  \item $\lnot\bigl(\lnot(P\land Q)\bigr)\equiv P\land Q$.
  \item $\lnot(P\Leftrightarrow Q)
        \equiv(P\land\lnot Q)\lor(Q\land\lnot P)$ (ou exclusif).
\end{enumerate}

%────────────────────────────────────────────────────────────────
\solutiontitle{Table \`a trois variables}

\renewcommand{\arraystretch}{1.2}
\begin{center}
\small
\begin{tabular}{|ccc|c|c|c|c|}
\hline
\rowcolor{BN!12}
$P$&$Q$&$R$ & $P\Rightarrow Q$ & $Q\Rightarrow R$ &
$(P\Rightarrow Q)\land(Q\Rightarrow R)$ & $P\Rightarrow R$ \\
\hline
\vrai&\vrai&\vrai & \vrai & \vrai & \vrai & \vrai \\
\vrai&\vrai&\faux & \vrai & \faux & \faux & \faux \\
\vrai&\faux&\vrai & \faux & \vrai & \faux & \vrai \\
\vrai&\faux&\faux & \faux & \vrai & \faux & \faux \\
\faux&\vrai&\vrai & \vrai & \vrai & \vrai & \vrai \\
\faux&\vrai&\faux & \vrai & \faux & \faux & \vrai \\
\faux&\faux&\vrai & \vrai & \vrai & \vrai & \vrai \\
\faux&\faux&\faux & \vrai & \vrai & \vrai & \vrai \\
\hline
\end{tabular}
\end{center}
Quand $(P\Rightarrow Q)\land(Q\Rightarrow R)$ est vraie, $P\Rightarrow R$ l'est aussi.
La r\'eciproque est fausse (ligne 6). Cela confirme la \textbf{transitivit\'e de l'implication}~:
si $P\Rightarrow Q$ et $Q\Rightarrow R$, alors $P\Rightarrow R$.\hfill$\square$

%────────────────────────────────────────────────────────────────
\solutiontitle{Tautologies et contradictions}

\begin{enumerate}
  \item $P\lor\lnot P$~: $P=\vrai\Rightarrow\vrai\lor\faux=\vrai$~;
        $P=\faux\Rightarrow\faux\lor\vrai=\vrai$. \textbf{Tautologie.}
  \item $P\land\lnot P$~: toujours \faux. \textbf{Contradiction.}
  \item $P\Rightarrow(Q\Rightarrow P)$~: si $P=\faux$, c'est vrai.
        Si $P=\vrai$, alors $Q\Rightarrow P=Q\Rightarrow\vrai=\vrai$.
        \textbf{Tautologie.}
  \item $(P\Rightarrow Q)\Leftrightarrow(\lnot P\lor Q)$~:
        $P\Rightarrow Q$ est vraie ssi $P=\faux$ ou $Q=\vrai$,
        exactement quand $\lnot P\lor Q$ est vraie. \textbf{Tautologie.}
\end{enumerate}

%────────────────────────────────────────────────────────────────
\solutiontitle{Divisibilit\'e et implication}

$P$~: $4\mid n$~; $Q$~: $2\mid n$.

\begin{enumerate}
  \item $P\Rightarrow Q$~: si $n=4k$ alors $n=2(2k)$~: \textbf{vraie.}
        $Q\Rightarrow P$~: fausse. Contre-exemple~: $n=2$, $2\mid2$ mais $4\nmid2$.
  \item $P\Leftrightarrow Q$ est fausse (les sens ne sont pas \'equivalents).
  \item Contrapos\'ee de $P\Rightarrow Q$~: si $2\nmid n$ alors $4\nmid n$.
        V\'erification directe~: si $n$ est impair, alors $n\ne4k$ pour tout entier $k$.
  \item Si $a\mid b$ et $b\mid c$, alors $a\mid c$ (\textbf{transitivit\'e}).
        Preuve~: $b=ka$ et $c=\ell b$, donc $c=\ell ka=(k\ell)a$, d'o\`u $a\mid c$.\hfill$\square$
\end{enumerate}

%────────────────────────────────────────────────────────────────
\solutiontitle{Quantificateurs imbriqu\'es}

\begin{enumerate}
  \item \og Pour tout r\'eel $x$, il existe un r\'eel $y$ tel que $y=x+1$.\fg\enspace
        \textbf{Vraie.} Prendre $y=x+1$.
        N\'egation~: $\exists x\in\R,\ \forall y\in\R,\ y\ne x+1$.
  \item \og Il existe un r\'eel $x$ tel que pour tout r\'eel $y$, $y=x+1$.\fg\enspace
        \textbf{Fausse.} $x+1$ ne peut pas \'egaler tous les r\'eels simultanément.
        N\'egation~: $\forall x\in\R,\ \exists y\in\R,\ y\ne x+1$. (Vraie.)
  \item \og Pour tout r\'eel $x$, il existe $y$ tel que $x+y=0$.\fg\enspace
        \textbf{Vraie.} Prendre $y=-x$.
        N\'egation~: $\exists x,\ \forall y,\ x+y\ne0$.
  \item \og Il existe un r\'eel $x$ tel que pour tout r\'eel $y$, $x\cdot y=y$.\fg\enspace
        \textbf{Vraie.} $x=1$ convient~: $1\cdot y=y$ pour tout $y\in\R$.
        N\'egation~: $\forall x\in\R,\ \exists y\in\R,\ x\cdot y\ne y$. (Fausse.)
\end{enumerate}

%────────────────────────────────────────────────────────────────
\solutiontitle{Raisonnement par contrapos\'ee}

\textbf{(1)} Contrapos\'ee~: si $n$ est pair, alors $n^2$ est pair.
$n=2k\Rightarrow n^2=4k^2=2(2k^2)$.\hfill$\square$

\textbf{(2)} Contrapos\'ee~: si $n$ est impair, alors $n^2$ est impair.
$n=2k+1\Rightarrow n^2=4k^2+4k+1=2(2k^2+2k)+1$.\hfill$\square$

\textbf{(3)} Contrapos\'ee~: si $n$ est pair, alors $n^2+2n$ est pair.
$n=2k\Rightarrow n^2+2n=4k^2+4k=2(2k^2+2k)$~: pair.\hfill$\square$
(L'hypoth\`ese peut \^etre vraie~: pour $n=1$, $n^2+2n=3$ est impair.)

\textbf{(4)} Contrapos\'ee de $ab\ne0\Rightarrow(a\ne0$ et $b\ne0)$~:
si $a=0$ ou $b=0$, alors $ab=0$. Trivial.\hfill$\square$

%────────────────────────────────────────────────────────────────
\solutiontitle{G\'eom\'etrie et logique}

$P$~: \og$ABC$ isoc\`ele en $A$\fg~; $Q$~: \og$\hat B=\hat C$\fg.

\begin{enumerate}
  \item $P\Rightarrow Q$~: si $AB=AC$, alors $\hat B=\hat C$ (propri\'et\'e du triangle isoc\`ele).
        \textbf{Vraie.}
  \item $Q\Rightarrow P$~: \og si $\hat B=\hat C$, alors $ABC$ est isoc\`ele en $A$\fg.
        \textbf{Vraie.} Preuve~: la somme des angles vaut $180^\circ$, donc deux angles
        \'egaux $\hat B=\hat C$ sont aigus~; le pied $H$ de la hauteur issue de $A$
        est alors sur le segment $[BC]$. Dans les triangles $ABH$ et $ACH$, rectangles en $H$~:
        $AH=AB\sin\hat B$ et $AH=AC\sin\hat C$. Comme $\sin\hat B=\sin\hat C$, on obtient
        $AB=AC$. (On peut aussi dire que $ABH$ et $ACH$ ont un c\^ot\'e commun $AH$ et les
        m\^emes angles, donc sont \'egaux.)
        Donc $P\Leftrightarrow Q$.
  \item \textbf{Contrapos\'ee de $P\Rightarrow Q$~:}
        $\lnot Q\Rightarrow\lnot P$, soit~:
        \og Si $\hat B\ne\hat C$, alors le triangle n'est pas isoc\`ele en $A$\fg.
        D\'emonstration~: si $\hat B\ne\hat C$, supposons par l'absurde $AB=AC$.
        Alors les angles de la base seraient \'egaux (propri\'et\'e de l'isoc\`ele),
        contradiction. Donc $AB\ne AC$~: le triangle n'est pas isoc\`ele en $A$.\hfill$\square$
  \item $\lnot P$~: \og $AB\ne AC$\fg\ (triangle non isoc\`ele en $A$).
        $\lnot Q$~: \og $\hat B\ne\hat C$\fg.
\end{enumerate}

%────────────────────────────────────────────────────────────────
\solutiontitle{N\'egation de propositions complexes}

\begin{enumerate}
  \item N\'egation~: \og Il existe un entier pair qui n'est ni divisible par $4$
        ni divisible par $6$.\fg
        En symboles~: $\exists n\in\Z,\ (2\mid n)\ \land\ (4\nmid n)\ \land\ (6\nmid n)$.
        (Vraie~: $n=2$ est pair, non divisible par $4$ et non divisible par $6$.)
  \item N\'egation~: \og Pour tout entier relatif $n$, $n^2\geq n$ ou $n\geq n^3$.\fg
        En symboles~: $\forall n\in\Z,\ n^2\geq n\ \lor\ n\geq n^3$.
        (Vraie~: $n^2-n=n(n-1)\geq0$ pour tout entier $n$, donc $n^2\geq n$.)
  \item N\'egation~: \og Il existe un r\'eel $x>0$ tel que, pour tout $y>0$, $y\geq x$.\fg
        En symboles~: $\exists x>0,\ \forall y>0,\ y\geq x$.
  \item N\'egation~: \og Il existe un nombre premier qui n'est ni impair ni
        \`egal \`a $2$.\fg\ Cette proposition est fausse~: le seul nombre premier
        pair est $2$.
\end{enumerate}

%────────────────────────────────────────────────────────────────
\solutiontitle{Erreurs classiques de raisonnement}

\begin{enumerate}
  \item \textbf{Erreur~: moyen terme non distribu\'e (faux syllogisme).}
        Les chats et les chiens sont des animaux, mais cela ne les rend pas identiques.
        La forme correcte serait~: si $P\Rightarrow R$ et $Q\Rightarrow R$, on ne peut pas
        conclure $P\Rightarrow Q$.
  \item \textbf{Erreur~: l'implication $x^2=4\Rightarrow x=2$ est fausse.}
        $x^2=4$ admet deux solutions~: $x=2$ et $x=-2$. On ne peut pas \'ecrire
        \og si $x^2=4$ alors $x=2$\fg.
  \item \textbf{Erreur~: g\'en\'eralisation h\^ative.}
        La v\'erification d'un nombre fini de cas ne prouve pas une assertion universelle.
        Avec $n=40$, on obtient $40^2+40+41=41^2$, qui n'est pas premier.
  \item \textbf{Erreur~: division par z\'ero.}
        Puisque $a=b$, on a $a-b=0$~: le passage de
        $(a-b)(a+b)=b(a-b)$ \`a $a+b=b$ est interdit.
\end{enumerate}

%────────────────────────────────────────────────────────────────
\solutiontitle{Double implication~: $n$ pair $\Leftrightarrow$ $n^2$ pair}

\textbf{Sens direct ($\Rightarrow$)~:} Si $n=2k$, alors $n^2=4k^2=2(2k^2)$~: pair.\hfill$\square$

\textbf{Sens r\'eciproque ($\Leftarrow$)~:} Par contrapos\'ee~: si $n$ est impair,
$n=2k+1$, donc $n^2=4k^2+4k+1=2(2k^2+2k)+1$~: impair.\hfill$\square$

Les deux sens \'etant d\'emontr\'es~: $n$ pair $\Leftrightarrow$ $n^2$ pair.

%────────────────────────────────────────────────────────────────
\solutiontitle{Conditions n\'ecessaires et suffisantes}

\begin{enumerate}
  \item $6\mid n$ pour $2\mid n$~: \textbf{condition suffisante.}
        $6\mid n\Rightarrow2\mid n$ (car $6=2\times3$).
        R\'eciproque fausse~: $n=2$, $2\mid2$ mais $6\nmid2$.
  \item $x=0$ pour $xy=0$~: \textbf{condition suffisante.}
        $x=0\Rightarrow xy=0$. R\'eciproque fausse~: $y=0$ suffit aussi.
  \item Un carr\'e pour un rectangle~: \textbf{condition suffisante},
        mais pas n\'ecessaire~: un rectangle de c\^ot\'es in\'egaux n'est pas un carr\'e.
  \item $\Delta=0$ pour une racine double~: \textbf{condition n\'ecessaire
        et suffisante} pour un trin\^ome du second degr\'e ($a\ne0$).
  \item $n$ impair pour $n^2$ impair~: \textbf{condition n\'ecessaire et suffisante}.
        Le sens direct est imm\'ediat~; si $n$ est pair, son carr\'e l'est aussi.
\end{enumerate}

%────────────────────────────────────────────────────────────────
\solutiontitle{D\'emonstration directe}

\begin{enumerate}
  \item $n=2k+1\Rightarrow n^3=(2k+1)^3=8k^3+12k^2+6k+1=2(4k^3+6k^2+3k)+1$~: impair.\hfill$\square$
  \item $n=2k+1\Rightarrow n^2+n=(2k+1)^2+(2k+1)=4k^2+4k+1+2k+1=4k^2+6k+2=2(2k^2+3k+1)$~: pair.\hfill$\square$
  \item $a=2p+1$, $b=2q+1\Rightarrow a+b=2p+2q+2=2(p+q+1)$~: pair.\hfill$\square$
  \item $a\mid b$ et $b\mid c$ signifient qu'il existe $k,\ell\in\Z$ tels que
        $b=ka$ et $c=\ell b$. Alors $c=\ell ka=(\ell k)a$,
        avec $\ell k\in\Z$. Ainsi $a\mid c$.\hfill$\square$
\end{enumerate}

%────────────────────────────────────────────────────────────────
\solutiontitle{Trouver l'erreur dans une preuve}

\begin{enumerate}
  \item \textbf{Erreur~: division par $n$ invalide pour $n=0$.}
        La division par $n$ n'est licite que si $n\ne0$.
        Or $n=0$ est un entier~: $0^2=0\geq0$ est vrai,
        mais l'\'etape \og diviser par $n$\fg\ est interdite pour $n=0$
        et ne prouve rien pour les autres valeurs.
        De plus, diviser par un entier $n<0$ \emph{renverse} l'in\'egalit\'e
        ($n^2\geq n$ donne alors $n\leq1$), et la conclusion \og $n\geq1$\fg\ est fausse
        pour tout $n\leq0$~: la \og preuve\fg\ aboutit \`a une affirmation fausse.
        Correction~: $n^2\geq n\Leftrightarrow n^2-n\geq0\Leftrightarrow n(n-1)\geq0$.
        Vrai si $n\leq0$ ou $n\geq1$, donc vrai pour tout entier.
        Mais faux pour les r\'eels~: $n=1/2$ donne $1/4<1/2$.
        La proposition est vraie pour les entiers, mais la \emph{preuve} reste invalide.
  \item \textbf{Erreur~: l'implication \og carr\'e des membres\fg\ n'est pas une \'equivalence.}
        De $a^2+b^2=(a+b)^2$ on d\'eduit $2ab=0$, donc $a=0$ ou $b=0$.
        Cela prouve~: $\sqrt{a^2+b^2}=a+b$ implique $a=0$ ou $b=0$.
        Mais la conclusion de la \og preuve\fg\ est invers\'ee~:
        elle affirme que $\sqrt{a^2+b^2}=a+b$ \emph{en g\'en\'eral}.
        Contre-exemple~: $a=3, b=4$~: $\sqrt{9+16}=5\ne7$.
\end{enumerate}

%────────────────────────────────────────────────────────────────
\solutiontitle{Irrationnalit\'e de $\sqrt{2}$}

\textbf{Th\'eor\`eme.} $\sqrt{2}\notin\Q$.

\textbf{D\'emonstration} (par l'absurde). Supposons $\sqrt{2}=p/q$ avec $\pgcd(p,q)=1$.
Alors $p^2=2q^2$, donc $2\mid p^2$, et par l'exercice~\ref{exo:1.21}(2)~: $2\mid p$.
Posons $p=2m$. Alors $4m^2=2q^2$, d'o\`u $q^2=2m^2$, donc $2\mid q$.
Mais $2\mid\pgcd(p,q)$~: contradiction.\hfill$\square$

\textit{G\'en\'eralisation.} Soit $k\geq2$ un entier qui n'est pas un carr\'e parfait.
Dans la d\'ecomposition de $k$ en facteurs premiers, les exposants ne sont pas tous pairs
(sinon $k$ serait un carr\'e)~: il existe un nombre premier $p$ dont l'exposant dans $k$ est \emph{impair}.
Supposons $\sqrt{k}=a/b$ avec $a,b\in\N^*$. Alors $a^2=kb^2$.
L'exposant de $p$ dans $a^2$ est pair (le double de celui de $a$).
Dans $kb^2$, c'est (exposant impair dans $k$) $+$ (exposant pair dans $b^2$), donc impair.
Par unicit\'e de la d\'ecomposition en facteurs premiers, c'est une contradiction~:
$\sqrt{k}\notin\Q$.\hfill$\square$

\textit{Attention.} Il ne suffit pas de prendre n'importe quel diviseur premier de $k$~:
pour $k=12=2^2\cdot3$ et $p=2$ (exposant pair), il n'y a pas de contradiction
(d'ailleurs $\sqrt{12}=2\sqrt{3}$). Il faut un premier d'exposant impair, ici $p=3$.

%────────────────────────────────────────────────────────────────
\solutiontitle{Irrationnalit\'e de $\sqrt{3}$}

\textbf{Lemme.} $3\mid n^2\Rightarrow3\mid n$.
\emph{Preuve.} Par contrapos\'ee~: si $3\nmid n$, alors $n=3k+1$ ou $n=3k+2$.
Si $n=3k+1$~: $n^2=9k^2+6k+1=3(3k^2+2k)+1$.
Si $n=3k+2$~: $n^2=9k^2+12k+4=3(3k^2+4k+1)+1$.
Dans les deux cas le reste de la division de $n^2$ par $3$ est $1$~: $3\nmid n^2$.\hfill$\square$

\textbf{Th\'eor\`eme.} Supposons $\sqrt{3}=p/q$ avec $\pgcd(p,q)=1$.
Alors $p^2=3q^2$, donc $3\mid p^2$, et par le lemme $3\mid p$~: $p=3m$.
Alors $9m^2=3q^2$, d'o\`u $3\mid q$~: contradiction.\hfill$\square$

%────────────────────────────────────────────────────────────────
\solutiontitle{Suite et logique}

\begin{enumerate}
  \item $n=0$~: $0$~; $n=1$~: $0$~; $n=2$~: $2$~; $n=3$~: $6$~; $n=4$~: $12$~: tous pairs.\checkmark
  \item $n^2-n=n(n-1)$. Deux entiers cons\'ecutifs sont toujours l'un pair et l'autre impair,
        donc leur produit est pair.\hfill$\square$
  \item Si $n=2k$~: $n^2=4k^2$, soit la forme $4k'$ avec $k'=k^2$.
        Si $n=2k+1$~: $n^2=4k^2+4k+1=4k(k+1)+1$, soit la forme $4k''+1$ avec $k''=k(k+1)$.
        Donc $n^2$ est de la forme $4k$ ou $4k+1$ pour tout $n\in\Z$.\hfill$\square$
\end{enumerate}

%────────────────────────────────────────────────────────────────
\solutiontitle{Infinit\'e des nombres premiers}

\textbf{Th\'eor\`eme} (Euclide). Il existe une infinit\'e de nombres premiers.

\textbf{D\'emonstration} (par l'absurde). Supposons qu'il n'en existe qu'un nombre fini
$p_1,\ldots,p_k$. Posons $N=p_1p_2\cdots p_k+1\geq2$.
$N$ admet un diviseur premier $p$. Mais pour tout $i$~:
$N = p_1p_2\cdots p_k+1$, donc $p_i\nmid N$ (reste $1$ dans la division)~: $p\notin\{p_1,\ldots,p_k\}$.
Contradiction.\hfill$\square$

%────────────────────────────────────────────────────────────────
\solutiontitle{Combinaisons logiques}

\begin{enumerate}
  \item \textbf{Exportation.} Table de v\'erit\'e~:
        \begin{center}\renewcommand{\arraystretch}{1.15}\small
        \begin{tabular}{|ccc|c|c|c|c|}
        \hline\rowcolor{BN!12}
        $P$&$Q$&$R$ & $Q\Rightarrow R$ & $P\Rightarrow(Q\Rightarrow R)$ & $P\land Q$ & $(P\land Q)\Rightarrow R$\\\hline
        \vrai&\vrai&\vrai & \vrai & \vrai & \vrai & \vrai\\
        \vrai&\vrai&\faux & \faux & \faux & \vrai & \faux\\
        \vrai&\faux&\vrai & \vrai & \vrai & \faux & \vrai\\
        \vrai&\faux&\faux & \vrai & \vrai & \faux & \vrai\\
        \faux&\vrai&\vrai & \vrai & \vrai & \faux & \vrai\\
        \faux&\vrai&\faux & \faux & \vrai & \faux & \vrai\\
        \faux&\faux&\vrai & \vrai & \vrai & \faux & \vrai\\
        \faux&\faux&\faux & \vrai & \vrai & \faux & \vrai\\\hline
        \end{tabular}
        \end{center}
        Les colonnes 5 et 7 co\"incident~: les deux formules ne sont fausses que pour
        $P=Q=\vrai$, $R=\faux$.\hfill$\square$
  \item $(P\Rightarrow Q)\land(P\Rightarrow R)\equiv P\Rightarrow(Q\land R)$.
        Directement~: le membre de gauche est faux exactement quand $P\Rightarrow Q$ ou
        $P\Rightarrow R$ est fausse, c'est-\`a-dire quand $P$ est vraie et ($Q$ fausse ou $R$ fausse).
        Le membre de droite est faux exactement quand $P$ est vraie et $Q\land R$ fausse,
        c'est-\`a-dire quand $P$ est vraie et ($Q$ fausse ou $R$ fausse).
        Les deux formules sont fausses sur les m\^emes lignes, donc \'equivalentes.\hfill$\square$
  \item $P\lor(Q\land R)\equiv(P\lor Q)\land(P\lor R)$~:
        distribution de $\lor$ sur $\land$. Par cas~:
        si $P=\vrai$, les deux c\^ot\'es sont vrais~;
        si $P=\faux$, gauche $=Q\land R$ et droite $=Q\land R$.
        Les deux cas couvrent les $8$ lignes de la table.\hfill$\square$
\end{enumerate}

%────────────────────────────────────────────────────────────────
\solutiontitle{Raisonner sur les r\'eels}

\begin{enumerate}
  \item On raisonne par cas, avec la r\`egle des signes d'un produit.
        Si $x\geq0$, alors $x^2=x\times x$ est un produit de deux nombres positifs, donc $x^2\geq0$.
        Si $x<0$, alors $x^2=x\times x$ est un produit de deux nombres n\'egatifs, donc $x^2>0$.
        Dans tous les cas $x^2\geq0$.
        Donc $x^2+1\geq0+1=1>0$.\hfill$\square$
  \item Par l'absurde~: supposons qu'il existe un plus petit r\'eel $\varepsilon>0$.
        Alors $\varepsilon/2>0$ et $\varepsilon/2<\varepsilon$~: contradiction.\hfill$\square$
  \item $0\leq(x-y)^2=x^2-2xy+y^2$, donc $2xy\leq x^2+y^2$~:
        $xy\leq\dfrac{x^2+y^2}{2}$.\hfill$\square$
\end{enumerate}

%────────────────────────────────────────────────────────────────
\solutiontitle{Nombres rationnels et puissances}

\begin{enumerate}
  \item $2^n$ est pair (car $2\mid2^n$) et $3^m$ est impair (car $3$ est impair,
        et une puissance d'impair est impaire). Un nombre ne peut \^etre
        \`a la fois pair et impair~: $2^n\ne3^m$.\hfill$\square$
  \item Pour $n\geq1$, $2^n$ et $3^n$ sont des produits de $n$ facteurs positifs,
        chaque facteur $2$ \'etant inf\'erieur au facteur $3$ correspondant~: $0<2^n<3^n$.
        Donc $0<2^n/3^n<1$. Aucun entier n'est strictement compris entre $0$ et $1$~:
        $2^n/3^n$ n'est jamais un entier.\hfill$\square$
  \item Valeurs approch\'ees~: $(2/3)^1\approx0{,}667$~; $(2/3)^5=32/243\approx0{,}132$~;
        $(2/3)^{10}\approx0{,}0173$~; $(2/3)^{20}\approx0{,}000\,301$.
        Passer de $n$ \`a $n+1$ revient \`a multiplier par $2/3<1$~: les valeurs diminuent,
        et elles semblent se rapprocher de $0$ (la notion de limite sera \'etudi\'ee plus tard).
        Mais $(2/3)^n$ n'est \textbf{jamais nulle}~: c'est un quotient de deux nombres
        strictement positifs, $2^n>0$ et $3^n>0$.
\end{enumerate}

%────────────────────────────────────────────────────────────────
\solutiontitle{R\'eciproque du th\'eor\`eme de Pythagore}

\begin{enumerate}
  \item \textbf{Contrapos\'ee~:} si $AB^2\ne AC^2+BC^2$, alors le triangle n'est pas
        rectangle en $C$.
  \item \textbf{R\'eciproque~:} si $AB^2=AC^2+BC^2$, alors le triangle est rectangle en $C$.
        Cette r\'eciproque est \textbf{vraie}~: c'est la r\'eciproque du th\'eor\`eme de Pythagore,
        un th\'eor\`eme du cours (admis ici, comme le demande la question~3).
  \item Le plus grand c\^ot\'e est $AB$. $5^2+12^2=25+144=169=13^2$~: par la r\'eciproque,
        le triangle de c\^ot\'es $5$, $12$, $13$ est rectangle (en $C$, sommet oppos\'e au c\^ot\'e $13$).
        De m\^eme, $6^2+8^2=36+64=100=10^2$~: celui de c\^ot\'es $6$, $8$, $10$ est rectangle.
  \item Contrapos\'ee de la r\'eciproque~: si le triangle n'est pas rectangle en $C$,
        alors $AB^2\ne AC^2+BC^2$. \textbf{Vraie}~: une implication et sa contrapos\'ee
        ont toujours la m\^eme valeur de v\'erit\'e, et la r\'eciproque est vraie.
\end{enumerate}

%────────────────────────────────────────────────────────────────
\subsection*{Probl\`emes}
\addcontentsline{toc}{subsection}{Probl\`emes}
%────────────────────────────────────────────────────────────────

\psolutiontitle{Circuit logique}

\renewcommand{\arraystretch}{1.25}
\begin{center}
\small
\begin{tabular}{|ccc|c|c|c|c|c|}
\hline
\rowcolor{BN!12}
$A$ & $B$ & $C$ & $B\lor C$ & $\lnot A$
  & $A\land(B\lor C)$ & $\lnot A\land B\land C$ & $\Phi$ \\
\hline
\vrai&\vrai&\vrai & \vrai & \faux & \vrai & \faux & \vrai \\
\vrai&\vrai&\faux & \vrai & \faux & \vrai & \faux & \vrai \\
\vrai&\faux&\vrai & \vrai & \faux & \vrai & \faux & \vrai \\
\vrai&\faux&\faux & \faux & \faux & \faux & \faux & \faux \\
\faux&\vrai&\vrai & \vrai & \vrai & \faux & \vrai & \vrai \\
\faux&\vrai&\faux & \vrai & \vrai & \faux & \faux & \faux \\
\faux&\faux&\vrai & \vrai & \vrai & \faux & \faux & \faux \\
\faux&\faux&\faux & \faux & \vrai & \faux & \faux & \faux \\
\hline
\end{tabular}
\end{center}

Le courant passe pour $(\vrai,\vrai,\vrai)$, $(\vrai,\vrai,\faux)$,
$(\vrai,\faux,\vrai)$, $(\faux,\vrai,\vrai)$~: exactement les cas o\`u au moins
deux interrupteurs sont ferm\'es. $\Phi\equiv(A\land B)\lor(A\land C)\lor(B\land C)$~:
c'est la \textbf{fonction majorit\'e}.\hfill$\square$

\psolutiontitle{Chevaliers et coquins}

$P$~: \og$X$ est chevalier\fg~; $Q$~: \og$Y$ est chevalier\fg~; $R$~: \og$Z$ est chevalier\fg.

\textbf{1. Traduction.} Un habitant est chevalier exactement quand sa d\'eclaration est vraie.
$X$ d\'eclare \og$Y$ est un coquin\fg, soit $\lnot Q$~: on a donc $P\Leftrightarrow\lnot Q$.
$Y$ d\'eclare \og$Z$ et moi sommes tous les deux des coquins\fg, soit $\lnot Q\land\lnot R$~:
on a donc $Q\Leftrightarrow(\lnot Q\land\lnot R)$.

\textbf{2. Les quatre cas.}
\begin{itemize}
  \item $P=\vrai$, $Q=\vrai$~: $X$ dit vrai, donc $Y$ est coquin~: contradiction avec $Q=\vrai$.
  \item $P=\faux$, $Q=\faux$~: $X$ ment, donc $Y$ n'est pas coquin~: contradiction avec $Q=\faux$.
  \item $P=\faux$, $Q=\vrai$~: $Y$ est chevalier, donc sa phrase est vraie, donc $Y$ est un coquin~:
        contradiction.
  \item $P=\vrai$, $Q=\faux$~: $X$ dit vrai ($Y$ est bien coquin), coh\'erent.
        $Y$ ment, donc \og $Z$ et $Y$ sont tous les deux coquins\fg\ est faux~;
        comme $Y$ est coquin, c'est que $Z$ n'est pas coquin~: $R=\vrai$.
\end{itemize}

\textbf{3. Conclusion.} Une seule possibilit\'e~: \textbf{$X$ est chevalier, $Y$ est coquin,
$Z$ est chevalier} (v\'erifi\'e sur les $8$ cas $(P,Q,R)$).
Remarque~: un chevalier ne peut jamais dire \og je suis un coquin\fg\ (ni \og moi et un autre
sommes coquins\fg)~; c'est ce qui \'elimine le troisi\`eme cas.

\psolutiontitle{\'Enigme logique~: qui a fait quoi~?}

Exactement une d\'eclaration est vraie. Testons les trois cas~:

\textbf{Cas $A$ coupable~:} $A$ dit \og ce n'est pas moi\fg~: \textbf{faux}.
$B$ dit \og c'est $A$\fg~: \textbf{vrai}.
$C$ dit \og ce n'est pas $A$\fg~: \textbf{faux}.
Une seule d\'eclaration vraie~: \textbf{coh\'erent.}

\textbf{Cas $B$ coupable~:} $A$ dit \og ce n'est pas moi\fg~: vrai.
$B$ dit \og c'est $A$\fg~: faux.
$C$ dit \og ce n'est pas $A$\fg~: vrai.
Deux d\'eclarations vraies~: \textbf{impossible.}

\textbf{Cas $C$ coupable~:} $A$~: vrai~; $B$~: faux~; $C$~: vrai.
Deux d\'eclarations vraies~: \textbf{impossible.}

\textbf{Conclusion.} \textbf{$A$ a commis le vol.}

\psolutiontitle{Le berger, le loup, la ch\`evre et le chou}

\'Etat~: $(B,L,C,X)$ o\`u $B$~: \og le berger\fg, $L$~: \og le loup\fg,
$C$~: \og la ch\`evre\fg, $X$~: \og le chou\fg, et $\vrai$ = c\^ot\'e d'arriv\'ee.
\'Etat initial~: $(\faux,\faux,\faux,\faux)$.
\'Etat final~: $(\vrai,\vrai,\vrai,\vrai)$.
Contraintes~: ne jamais avoir $L=C\ne B$ (loup mange ch\`evre)
ni $C=X\ne B$ (ch\`evre mange chou).

\textbf{\'Etats interdits.} Ce sont les $6$ \'etats~:
$(\faux,\faux,\vrai,\vrai)$, $(\faux,\vrai,\vrai,\faux)$, $(\faux,\vrai,\vrai,\vrai)$,
$(\vrai,\faux,\faux,\faux)$, $(\vrai,\faux,\faux,\vrai)$, $(\vrai,\vrai,\faux,\faux)$
(le berger est de l'autre c\^ot\'e et laisse ensemble loup et ch\`evre, ou ch\`evre et chou).

\textbf{Solution~:}
\begin{enumerate}
  \item Traverser avec la ch\`evre~: $(\vrai,\faux,\vrai,\faux)$.
  \item Revenir seul~: $(\faux,\faux,\vrai,\faux)$.
  \item Traverser avec le loup~: $(\vrai,\vrai,\vrai,\faux)$.
  \item Revenir avec la ch\`evre~: $(\faux,\vrai,\faux,\faux)$.
  \item Traverser avec le chou~: $(\vrai,\vrai,\faux,\vrai)$.
  \item Revenir seul~: $(\faux,\vrai,\faux,\vrai)$.
  \item Traverser avec la ch\`evre~: $(\vrai,\vrai,\vrai,\vrai)$.
\end{enumerate}
Solution alternative~: \'echanger \og loup\fg\ et \og chou\fg\ aux \'etapes~3 et~5
(l'\'etape~4 ram\`ene la ch\`evre dans les deux cas).
Un parcours en largeur sur les \'etats autoris\'es montre qu'il n'existe pas de solution en moins de $7$ travers\'ees
et que ces deux solutions sont les seules en $7$ travers\'ees~; il y en a d'autres, plus longues,
si l'on accepte de refaire des allers-retours inutiles.

\psolutiontitle{Parit\'e des puissances}

\begin{enumerate}
  \item Si $n$ est pair~: $n=2k$, $n^3=8k^3=2(4k^3)$~: pair.
        Si $n$ est impair~: $n=2k+1$, $n^3=(2k+1)^3=2(\ldots)+1$~: impair.
        Donc $n$ et $n^3$ ont la m\^eme parit\'e.\hfill$\square$
  \item $n=2m+1$~: $n^2=(2m+1)^2=4m^2+4m+1=4m(m+1)+1$.
        Parmi $m$ et $m+1$, l'un est pair~: $m(m+1)=2k$.
        Donc $n^2=8k+1$.\hfill$\square$
  \item Si $4\mid(a^2+b^2)$, alors $a^2+b^2$ est de la forme $4\ell$.
        Or $a^2$ et $b^2$ sont chacun de la forme $4k$ ou $4k+1$
        (voir exercice~\ref{exo:1.31}). Leur somme vaut $4\ell$ seulement si
        les deux sont de la forme $4k$, c'est-\`a-dire $a$ et $b$ sont pairs.\hfill$\square$
\end{enumerate}

\psolutiontitle{Divisibilit\'e et parit\'e des carr\'es}

\begin{enumerate}
  \item $n=2k\Rightarrow n^2=4k^2$~: de la forme $4k'$.
        $n=2k+1\Rightarrow n^2=4k^2+4k+1=4k(k+1)+1$~: de la forme $4k'+1$.
        Donc $n^2$ est de la forme $4k$ ou $4k+1$.\hfill$\square$
  \item $a^2+b^2$ vaut $4k$, $4k+1$ ou $4k+2$ (jamais $4k+3$).
        Donc $a^2+b^2$ n'est jamais de la forme $4k+3$.
        Les entiers de cette forme sont~: $3,7,11,15,\ldots$~:
        aucun d'eux ne s'\'ecrit comme somme de deux carr\'es.\hfill$\square$
  \item Les carr\'es inf\'erieurs ou \'egaux \`a $8$ sont $0,1,4$ (on autorise $0$).
        Leurs sommes de deux~: $0,1,2,4,5,8$. Donc
        $1=1+0$, $2=1+1$, $4=4+0$, $5=4+1$, $8=4+4$ sont possibles,
        et les valeurs \textbf{impossibles} sont $3$, $6$ et $7$.
        Les entiers $3$ et $7$ sont de la forme $4k+3$ (question~2).
        En revanche $6=4\cdot1+2$ n'est pas de cette forme~: la question~2 ne l'exclut pas,
        et on le constate par la liste finie ci-dessus ($6$ n'est pas dans la liste des sommes).
  \item Sommes de deux carr\'es dans $\{1,\ldots,20\}$~:
        $1,2,4,5,8,9,10,13,16,17,18,20$. Les impossibles~: $3,6,7,11,12,14,15,19$.
\end{enumerate}

\psolutiontitle{Logique dans les preuves g\'eom\'etriques}

\begin{enumerate}
  \item $P\Rightarrow Q$ ($\vrai$)~; $P\Rightarrow R$ ($\vrai$)~;
        $Q\Rightarrow S$ ($\vrai$)~; $R\Rightarrow S$ ($\vrai$)~; $P\Rightarrow S$ ($\vrai$).
        R\'eciproquement~: $Q\not\Rightarrow P$ (un rectangle non carr\'e)~;
        $R\not\Rightarrow P$ (un losange non carr\'e).
        Donc $P\Rightarrow Q$ et $P\Rightarrow R$ mais $Q\not\Leftrightarrow P$.
  \item Contrapos\'ee de $P\Rightarrow Q$~: $\lnot Q\Rightarrow\lnot P$, c'est-\`a-dire
        \og si $ABCD$ n'est pas un rectangle, alors ce n'est pas un carr\'e\fg.
        Contrapos\'ee de $Q\Rightarrow S$~: $\lnot S\Rightarrow\lnot Q$, c'est-\`a-dire
        \og si $ABCD$ n'est pas un parall\'elogramme, alors ce n'est pas un rectangle\fg.
  \item On veut d\'emontrer $\lnot S\Rightarrow\lnot P$, qui est la contrapos\'ee de $P\Rightarrow S$.
        Il suffit donc de d\'emontrer $P\Rightarrow S$.
        Supposons que $ABCD$ est un carr\'e~; ses quatre angles sont droits.
        On utilise le th\'eor\`eme (du plan)~: \emph{si deux droites sont perpendiculaires \`a une
        m\^eme droite, alors elles sont parall\`eles entre elles.}
        Les angles en $B$ et en $C$ sont droits, donc $(AB)\perp(BC)$ et $(DC)\perp(BC)$~:
        $(AB)\parallel(DC)$.
        Les angles en $A$ et en $B$ sont droits, donc $(AD)\perp(AB)$ et $(BC)\perp(AB)$~:
        $(AD)\parallel(BC)$.
        Les c\^ot\'es oppos\'es de $ABCD$ sont parall\`eles deux \`a deux~:
        par d\'efinition, $ABCD$ est un parall\'elogramme.
        Par contrapos\'ee, si $ABCD$ n'est pas un parall\'elogramme, ce n'est pas un carr\'e.\hfill$\square$
        (Autre voie~: d'apr\`es la question~1, $P\Rightarrow Q$ et $Q\Rightarrow S$, donc
        $P\Rightarrow S$ par transitivit\'e~; la preuve ci-dessus d\'emontre en fait
        $Q\Rightarrow S$, car elle n'utilise que les angles droits.)
  \item N\'egation de $\forall\,ABCD\text{ carr\'e},\ (Q\land R)$~:
        $\exists\,ABCD\text{ carr\'e},\ \lnot(Q\land R)$, soit
        \og il existe un carr\'e qui n'est pas un rectangle ou qui n'est pas un losange\fg
        ($\lnot Q\lor\lnot R$, De Morgan).
        Cette n\'egation est \textbf{fausse}~: un carr\'e a quatre angles droits
        (c'est un rectangle) et quatre c\^ot\'es \'egaux (c'est un losange),
        donc aucun carr\'e n'est un contre-exemple. La proposition de d\'epart est vraie.\hfill$\square$
  \item[] \textit{Remarque.} La r\'eciproque est aussi vraie~: un rectangle qui est un losange
        est un carr\'e, donc $P\Leftrightarrow(Q\land R)$.
\end{enumerate}

\psolutiontitle{Sudoku $4\times4$ par d\'eduction logique}

Grille de d\'epart~:
\[
\begin{array}{|cc|cc|}
\hline
  & 2 & 3 &   \\
3 &   &   & 2 \\
\hline
  & 3 & 4 &   \\
4 &   &   & 3 \\
\hline
\end{array}
\]

\textbf{\'Etape~1.} Colonne~1 contient $\{3,4\}$~: manquent $1$ et $2$.
Ligne~1 contient $\{2,3\}$~: manquent $1$ et $4$.
$(1,1)$ doit \^etre dans $\{1,2\}\cap\{1,4\}=\{1\}$~: \textbf{$(1,1)=1$}.
Ligne~1 compl\`ete~: $(1,4)=4$.

\textbf{\'Etape~2.} Colonne~4 contient $\{4,2,3\}$~: il manque $1$, et la seule case vide
de la colonne~4 est $(3,4)$~: \textbf{$(3,4)=1$}.
Ligne~3 contient alors $\{3,4,1\}$~: il manque $2$, donc \textbf{$(3,1)=2$}.

\textbf{\'Etape~3.} Colonne~2 contient $\{2,3\}$~: manquent $1$ et $4$.
Ligne~2 contient $\{3,2\}$~: manquent $1$ et $4$.
Bloc haut-gauche contient $\{1,2,3\}$~: manque $4$.
$(2,2)\in\{1,4\}\cap\{1,4\}\cap\{4\}=\{4\}$~: \textbf{$(2,2)=4$}. Donc $(2,3)=1$.

\textbf{\'Etape~4.} Colonne~2 contient maintenant $\{2,4,3\}$~: il manque $1$,
donc \textbf{$(4,2)=1$}.
Ligne~4 contient $\{4,1,3\}$~: il manque $2$, donc \textbf{$(4,3)=2$}.

\textbf{Solution unique~:}
\[
\begin{array}{|cc|cc|}
\hline
\mathbf{1} & 2 & 3 & \mathbf{4} \\
3 & \mathbf{4} & \mathbf{1} & 2 \\
\hline
\mathbf{2} & 3 & 4 & \mathbf{1} \\
4 & \mathbf{1} & \mathbf{2} & 3 \\
\hline
\end{array}
\]

\begin{enumerate}
  \item Chaque \'etape utilise une implication du type~:
        \og La case $(i,j)$ ne peut pas contenir $k$ car $k$ est
        d\'ej\`a pr\'esent dans la ligne/colonne/bloc~: donc elle contient~$\ell$.\fg
  \item $8$ implications ont \'et\'e utilis\'ees, une par case vide
        (la grille de d\'epart a $8$ cases vides), regroup\'ees en $4$ \'etapes de $2$ cases.
  \item La solution est unique~: chaque \'etape \'etait forc\'ee.
\end{enumerate}

\psolutiontitle{Sommes et formules alg\'ebriques}

\begin{enumerate}
  \item V\'erification~: $n=1$~: $1=1$\checkmark~; $n=2$~: $3=3$\checkmark~;
        $n=3$~: $6=6$\checkmark~; $n=4$~: $10=10$\checkmark~; $n=5$~: $15=15$\checkmark.
  \item Deux m\'ethodes pour $n=5$~: $1+2+3+4+5=15$ et $5\cdot6/2=15$. \'Egaux.\checkmark
  \item \textbf{Algorithme~:}
\begin{algorithm}[H]
\caption{V\'erification de la formule de Gauss}
\KwIn{un entier $n\geq1$}
$S \leftarrow 0$\;
\Pour{$k\leftarrow1$ \KwTo $n$}{$S\leftarrow S+k$\;}
\Si{$S = n\times(n+1)/2$}{\Afficher \og Formule v\'erifi\'ee\fg\;}
\Sinon{\Afficher \og Erreur\fg\;}
\end{algorithm}
        Pour $n=10$~: $S=55=10\times11/2$. Pour $n=100$~: $S=5050=100\times101/2$.
        Dans les deux cas l'algorithme affiche \og Formule v\'erifi\'ee\fg.
  \item Tableau~: $n=1,2,3,4,5$ donne $1,\ 1+3=4,\ 1+3+5=9,\ 16,\ 25$.
        On conjecture $1+3+5+\cdots+(2n-1)=n^2$, et on le v\'erifie pour $n=1,\ldots,5$
        (valeurs ci-dessus).
        \emph{Pour aller plus loin (sans calcul).} Disposer des points en carr\'e~:
        un carr\'e de $1\times1$ point, puis on l'entoure d'une \og \'equerre\fg\ de $3$ points
        pour obtenir un carr\'e $2\times2$, puis d'une \'equerre de $5$ points pour un carr\'e
        $3\times3$, et ainsi de suite~: chaque nouveau nombre impair compl\`ete le carr\'e suivant,
        donc $1+3+\cdots+(2n-1)$ points forment un carr\'e $n\times n$.
        \emph{Preuve directe.} \'Ecrivons la somme $S$ dans les deux sens~:
        $S=1+3+\cdots+(2n-1)$ et $S=(2n-1)+(2n-3)+\cdots+1$.
        En additionnant terme \`a terme, chaque colonne vaut $2n$
        (le $k$-i\`eme terme $2k-1$ et le $k$-i\`eme terme $2n-2k+1$ ont pour somme $2n$),
        et il y a $n$ colonnes~: $2S=n\times2n$, donc $S=n^2$.
\end{enumerate}

\psolutiontitle{Codage et logique~: parit\'e}

\begin{enumerate}
  \item \texttt{1011}~: trois $1$~: nombre impair. Bit de parit\'e $= 1$
        (pour que le total soit pair~: $3+1=4$, pair).
        Message envoy\'e~: \texttt{10111}.
  \item \texttt{10110}~: trois $1$~: nombre impair. \textbf{Erreur d\'etect\'ee.}
  \item \texttt{10100}~: deux $1$~: nombre pair. Aucune erreur d\'etect\'ee
        (alors que ce mot diff\`ere de \texttt{10111} en deux positions~:
        c'est le cas de la question~4).
  \item Si deux erreurs se produisent, la parit\'e reste inchang\'ee~:
        le syst\`eme \textbf{ne peut pas} d\'etecter deux erreurs simultan\'ees.
        La parit\'e ne d\'etecte que les erreurs en nombre \emph{impair}.
\end{enumerate}

\psolutiontitle{D\'eduction logique~: tournoi de foot}

6 matchs au total (6 paires parmi 4 \'equipes). $A$ gagne 3 matchs~: 9 points.

$B$ a 6 points. Comme $B$ a perdu contre $A$, elle a gagn\'e ses deux autres
matchs. Ainsi $B$ bat $C$ et $D$. Les points de $C$ et $D$ proviennent alors
uniquement de leur confrontation directe. Pour qu'ils soient \`a \`egalit\'e,
ce match est nul.
(V\'erification~: $5$ matchs d\'ecisifs $\times3+1$ nul $\times2=17=9+6+1+1$.)

\textbf{Classement final~:} $A=9$, $B=6$, $C=D=1$.\hfill$\square$

\psolutiontitle{Logique et calcul alg\'ebrique}

\begin{enumerate}
  \item N\'egation~: \og Il existe $x\in\R$ tel que $x^2<0$.\fg
        Fausse~: la proposition originale est vraie.
  \item Contrapos\'ee~: \og Si $x\ne0$, alors $x^2\ne0$.\fg
        Preuve~: si $x\ne0$, alors $x^2=x\cdot x\ne0$ (produit de deux non-nuls).
  \item Pour $a\ne0$, $ax^2+bx+c=0$ a deux racines r\'eelles distinctes
        $\Leftrightarrow\Delta>0$.
        R\'eciproque~: $\Delta>0\Rightarrow$ deux racines distinctes. \textbf{Vraie.}
  \item \og$\exists x\in\R,\ x^2+1=0$\fg~: n\'egation~:
        \og$\forall x\in\R,\ x^2+1\ne0$.\fg La n\'egation est vraie
        car $x^2\geq0$, donc $x^2+1\geq1>0$.
\end{enumerate}

\psolutiontitle{Algorithme et divisibilit\'e}

\begin{enumerate}
  \item Oui~: $n\bmod3=0\Leftrightarrow3\mid n$ et $n\bmod9\ne0\Leftrightarrow9\nmid n$.
        Les deux notations sont \'equivalentes.
  \item N\'egation~: \og$3\nmid n$ ou $9\mid n$\fg.
  \item Pour $n=3$~: $3\bmod9=3\ne0$~: affich\'e.
        $n=6$~: affich\'e. $n=9$~: $9\bmod9=0$~: non affich\'e.
        $n=12$~: affich\'e. $n=18$~: non affich\'e. $n=21$~: affich\'e.
  \item Entre $1$ et $100$~: multiples de $3$~: $\lfloor100/3\rfloor=33$.
        Multiples de $9$~: $\lfloor100/9\rfloor=11$.
        Multiples de $3$ mais pas $9$~: $33-11=\mathbf{22}$ entiers.
\end{enumerate}

%────────────────────────────────────────────────────────────────
\subsection*{D\'efis}
\addcontentsline{toc}{subsection}{D\'efis}
%────────────────────────────────────────────────────────────────

\noindent\phantomsection
{\color{BO}\bfseries\small D\'efi~1~\textendash\ \textit{Le paradoxe du menteur}}
\par\smallskip

\begin{enumerate}
  \item Si elle est vraie~: elle affirme qu'elle est fausse, donc elle est fausse.
        Si elle est fausse~: elle affirme qu'elle est fausse, ce qui est vrai.
        Dans les deux cas on obtient une contradiction.
  \item Cette phrase n'est pas une proposition car elle ne peut avoir de valeur de
        v\'erit\'e d\'etermin\'ee~: elle est \textbf{auto-r\'ef\'erentielle}.
        Le cadre logique du cours exige que les propositions aient une valeur de
        v\'erit\'e bien d\'efinie.
  \item Autres exemples~: \og Cette phrase contient cinq mots.\fg\
        (Non paradoxale~: vraie.) \og Cette phrase est vraie.\fg\
        (Non paradoxale, mais elle n'est pas une proposition au sens du cours~:
        on peut lui donner la valeur \og vrai\fg\ ou la valeur \og faux\fg\ sans contradiction,
        donc rien ne fixe sa valeur de v\'erit\'e. Ce n'est pas une tautologie.)
        \og Cette phrase est fausse\fg\ est la version paradoxale.
  \item Soit $R=\{x\mid x\notin x\}$. Si $R\in R$, alors $R$ v\'erifie la condition de
        d\'efinition, donc $R\notin R$. Si $R\notin R$, alors $R$ v\'erifie la condition,
        donc $R\in R$. Dans les deux cas on obtient une contradiction~:
        $R\in R\Leftrightarrow R\notin R$. C'est la m\^eme structure que le menteur
        (\og elle est vraie $\Leftrightarrow$ elle est fausse\fg), avec l'appartenance \`a la place de la v\'erit\'e.
\end{enumerate}

\noindent\phantomsection
{\color{BO}\bfseries\small D\'efi~2~\textendash\ \textit{Construire sa propre table de v\'erit\'e}}
\par\smallskip

\begin{enumerate}
  \item Table du \textbf{ou exclusif} $P\oplus Q$~:
        \begin{center}\renewcommand{\arraystretch}{1.2}
        \begin{tabular}{|cc|c|}
        \hline\rowcolor{BN!10}$P$&$Q$&$P\oplus Q$\\\hline
        \vrai&\vrai&\faux\\
        \vrai&\faux&\vrai\\
        \faux&\vrai&\vrai\\
        \faux&\faux&\faux\\\hline
        \end{tabular}
        \end{center}
  \item $P\oplus Q\equiv(P\lor Q)\land\lnot(P\land Q)$~:
        v\'erifier ligne par ligne~: vrai seulement quand exactement un des deux est vrai.\checkmark
  \item \textbf{\'Etape~a~: le \og ou\fg\ avec $\Rightarrow$.}
        $\lnot P\Rightarrow Q$ est fausse seulement quand $\lnot P$ est vraie et $Q$ fausse,
        c'est-\`a-dire quand $P$ et $Q$ sont fausses~: exactement comme $P\lor Q$.
        \begin{center}\renewcommand{\arraystretch}{1.15}
        \begin{tabular}{|cc|c|c||c|c|c|}
        \hline\rowcolor{BN!10}
        $P$&$Q$&$P\lor Q$&$\lnot P\Rightarrow Q$&$\lnot Q$&$P\Rightarrow\lnot Q$&$\lnot(P\Rightarrow\lnot Q)$\\\hline
        \vrai&\vrai&\vrai&\vrai&\faux&\faux&\vrai\\
        \vrai&\faux&\vrai&\vrai&\vrai&\vrai&\faux\\
        \faux&\vrai&\vrai&\vrai&\faux&\vrai&\faux\\
        \faux&\faux&\faux&\faux&\vrai&\vrai&\faux\\\hline
        \end{tabular}
        \end{center}
        \textbf{\'Etape~b~: le \og et\fg\ avec $\Rightarrow$.}
        $P\Rightarrow\lnot Q$ est fausse seulement quand $P$ et $Q$ sont vraies~;
        sa n\'egation $\lnot(P\Rightarrow\lnot Q)$ est donc vraie seulement quand $P$ et $Q$
        sont vraies~: c'est $P\land Q$ (derni\`ere colonne de la table).
        Cette \'equivalence vaut pour deux propositions quelconques $A$ et $B$~:
        $A\land B\equiv\lnot(A\Rightarrow\lnot B)$.

        \textbf{\'Etape~c~: assemblage.} On part de $P\oplus Q\equiv(P\lor Q)\land\lnot(P\land Q)$
        (question~2) et on remplace \`a l'aide des \'etapes~a et~b~:
        \begin{itemize}
          \item $P\lor Q$ devient $\lnot P\Rightarrow Q$~;
          \item $\lnot(P\land Q)$ devient $\lnot\lnot(P\Rightarrow\lnot Q)$, c'est-\`a-dire $P\Rightarrow\lnot Q$~;
          \item le $\land$ central, appliqu\'e \`a $A=(\lnot P\Rightarrow Q)$ et $B=(P\Rightarrow\lnot Q)$,
                devient $\lnot(A\Rightarrow\lnot B)$ (\'etape~b).
        \end{itemize}
        D'o\`u
        \[P\oplus Q\ \equiv\ \lnot\bigl((\lnot P\Rightarrow Q)\Rightarrow\lnot(P\Rightarrow\lnot Q)\bigr),\]
        qui n'utilise que $\Rightarrow$ et $\lnot$.

        \textbf{V\'erification.} $\lnot P\Rightarrow Q$ est fausse seulement si $P=Q=\faux$~;
        $P\Rightarrow\lnot Q$ est fausse seulement si $P=Q=\vrai$.
        Pour $P\ne Q$, $A=B=\vrai$, donc $A\Rightarrow\lnot B$ est fausse et la formule est vraie~;
        pour $P=Q$, l'une des deux est fausse, $A\Rightarrow\lnot B$ est vraie et la formule est fausse.
        On retrouve la table de $P\oplus Q$.\checkmark

        \emph{Remarque.} Une formule plus courte existe~: $P\oplus Q\equiv\lnot(P\Leftrightarrow Q)$,
        et $\lnot\bigl((P\Rightarrow Q)\land(Q\Rightarrow P)\bigr)\equiv(P\Rightarrow Q)\Rightarrow\lnot(Q\Rightarrow P)$
        (\'etape~b), d'o\`u $P\oplus Q\equiv(P\Rightarrow Q)\Rightarrow\lnot(Q\Rightarrow P)$
        (v\'erifi\'e sur les $4$ lignes).
  \item $P\oplus P$ vaut $\faux$ (les deux valeurs sont toujours \'egales, donc jamais exactement une vraie)
        et $P\oplus\lnot P$ vaut $\vrai$ (exactement une des deux est vraie), pour toute valeur de $P$.\checkmark
        \emph{Remarque.} Avec seulement $\oplus$ et les constantes $\vrai$, $\faux$, on ne peut obtenir que
        des fonctions de la forme $c\oplus aP\oplus bQ$ (affines) : $P\land Q$ n'en fait pas partie.
        Pour obtenir toutes les fonctions il faut ajouter $\land$~:
        $\lnot P=P\oplus\vrai$ et $P\lor Q=P\oplus Q\oplus(P\land Q)$.
  \item Chiffrement~: $C\oplus K=(M\oplus K)\oplus K=M\oplus(K\oplus K)=M\oplus\faux=M$
        (associativit\'e de $\oplus$ et $K\oplus K=\faux$). Le d\'echiffrement est donc la m\^eme op\'eration.\checkmark
\end{enumerate}

\noindent\phantomsection
{\color{BO}\bfseries\small D\'efi~3~\textendash\ \textit{\'Enigmes logiques~: chevaliers et coquins}}
\par\smallskip

\begin{enumerate}
  \item $A$ dit \og nous sommes tous les deux des coquins\fg.
        Si $A$ est chevalier~: la phrase est vraie, donc $A$ est coquin~: contradiction.
        Donc \textbf{$A$ est coquin} (il ment). La phrase est fausse, donc
        ils ne sont \emph{pas} tous les deux coquins~: \textbf{$B$ est chevalier}.
  \item $C$ dit \og Au moins l'un d'entre nous est un coquin.\fg\ ; $D$ dit \og $C$ est un coquin.\fg
        Si $C$ est coquin~: sa phrase est fausse, donc aucun des trois n'est coquin~:
        contradiction. Donc \textbf{$C$ est chevalier}, sa phrase est vraie, et il y a un coquin.
        Alors $D$ dit \og $C$ est un coquin\fg, ce qui est faux~: \textbf{$D$ est coquin}.
        $E$ ne dit rien~: sa nature n'est pas d\'etermin\'ee.
        Il y a donc exactement \textbf{deux} possibilit\'es~:
        $(C,D,E)=(\text{chevalier},\text{coquin},\text{chevalier})$ ou
        $(\text{chevalier},\text{coquin},\text{coquin})$.
  \item Exemple~: $F$ dit \og $G$ est un coquin\fg, $G$ dit \og $F$ et $H$ sont du m\^eme type\fg,
        $H$ dit \og $F$ est un coquin\fg.
        Si $F$ est coquin, $G$ est chevalier, $H$ dit vrai donc est chevalier, et $G$ affirme que $F$, $H$
        sont du m\^eme type~: faux, contradiction. Donc $F$ est chevalier, $G$ est coquin~;
        $H$ dit \og $F$ est un coquin\fg, faux, donc $H$ est coquin~;
        et $G$ (coquin) dit \og $F$, $H$ m\^eme type\fg, ce qui est bien faux.
        Solution unique~: $F$ chevalier, $G$ et $H$ coquins (v\'erifi\'e sur les $8$ cas).
  \item $A$ est chevalier $\Leftrightarrow$ sa phrase est vraie~: $P\Leftrightarrow(\lnot P\land\lnot Q)$.
        \begin{center}\renewcommand{\arraystretch}{1.2}
        \begin{tabular}{|cc|c|c|}
        \hline\rowcolor{BN!10}$P$&$Q$&$\lnot P\land\lnot Q$&$P\Leftrightarrow(\lnot P\land\lnot Q)$\\\hline
        \vrai&\vrai&\faux&\faux\\
        \vrai&\faux&\faux&\faux\\
        \faux&\vrai&\faux&\vrai\\
        \faux&\faux&\vrai&\faux\\\hline
        \end{tabular}
        \end{center}
        Seule la ligne $P=\faux$, $Q=\vrai$ rend la formule vraie~: $A$ est coquin et $B$ est chevalier.
\end{enumerate}

\noindent\phantomsection
{\color{BO}\bfseries\small D\'efi~4~\textendash\ \textit{Bijections et taille d'ensembles}}
\par\smallskip

\begin{enumerate}
  \item $\{1,2,3\}\to\{a,b,c\}$~: bijection $1\mapsto a$, $2\mapsto b$, $3\mapsto c$. \textbf{Oui.}
  \item Bijection $\N\to\Z$~: $0\mapsto0$, $1\mapsto1$, $2\mapsto-1$, $3\mapsto2$, $4\mapsto-2$, \ldots
        En g\'en\'eral~: $f(2k)=-k$ et $f(2k+1)=k+1$ (pour $k\in\N$).
        Les entiers pairs donnent les entiers n\'egatifs ou nuls, les impairs les entiers strictement positifs,
        chaque entier de $\Z$ est atteint exactement une fois~: $f$ est bijective.
        Donc $\N$ et $\Z$ ont \og autant\fg\ d'\'el\'ements~: \textbf{\'equipotents.}
  \item $n\mapsto2n$ est une bijection de $\N$ sur l'ensemble des entiers pairs
        (injective, et tout entier pair $2m$ a pour ant\'ec\'edent $m$).
        Oui~: il y en a autant, bien que les pairs forment une partie stricte de $\N$ (Galil\'ee).
        \emph{Remarque.} De m\^eme $\N\times\N$ est en bijection avec $\N$ en parcourant les diagonales~:
        $(0,0),(1,0),(0,1),(2,0),(1,1),(0,2),\ldots$
  \item Raisonnons par l'absurde~: supposons qu'on puisse lister tous les r\'eels de $[0,1]$
        dans l'ordre croissant, $r_1<r_2<r_3<\cdots$, sans en oublier aucun, chacun \'etant
        suivi du \og suivant\fg.
        Le r\'eel $c=\dfrac{r_1+r_2}{2}$ est dans $[0,1]$ et v\'erifie $r_1<c<r_2$
        (car $c-r_1=\dfrac{r_2-r_1}{2}>0$ et $r_2-c=\dfrac{r_2-r_1}{2}>0$).
        Il devrait figurer dans la liste, mais il n'y a aucune place entre $r_1$ et $r_2$~:
        contradiction. Un r\'eel n'a donc pas de \og suivant\fg, contrairement \`a un entier
        ($n$ est suivi de $n+1$).
        \emph{Attention~:} cela montre seulement qu'on ne peut pas lister les r\'eels de $[0,1]$
        \emph{dans l'ordre}. Le r\'esultat de Cantor (aucune liste, dans aucun ordre, ne les
        contient tous) est plus fort et reste admis \`a ce niveau.
\end{enumerate}

\noindent\phantomsection
{\color{BO}\bfseries\small D\'efi~5~\textendash\ \textit{Raisonnement par l'absurde~: deux preuves classiques}}
\par\smallskip

\begin{enumerate}
  \item Supposition~: $7k$ est le plus grand multiple de $7$.
        Alors $7(k+1)=7k+7$ est aussi un multiple de $7$ et $7(k+1)>7k$~: contradiction.
        Donc il y a une infinit\'e de multiples de $7$.\hfill$\square$
  \item Supposition~: il existe un plus grand entier $N$.
        Alors $N+1$ est un entier et $N+1>N$~: contradiction.\hfill$\square$
  \item Supposons $\sqrt{5}=p/q$ avec $\pgcd(p,q)=1$. Alors $p^2=5q^2$, donc $5\mid p^2$.
        Si $5\nmid p$, alors $p\equiv1,2,3$ ou $4\pmod5$ et $p^2\equiv1,4,4$ ou $1\pmod5$, jamais $0$~;
        donc $5\mid p$. Posons $p=5m$~: $25m^2=5q^2$, d'o\`u $q^2=5m^2$ et $5\mid q^2$, donc $5\mid q$.
        Alors $5\mid\pgcd(p,q)=1$~: contradiction. Donc $\sqrt{5}\notin\Q$.\hfill$\square$
  \item Soient $a<b$ r\'eels et $c=\dfrac{a+b}{2}$. Alors $c-a=\dfrac{b-a}{2}>0$ et
        $b-c=\dfrac{b-a}{2}>0$, donc $a<c<b$. Il n'existe donc pas deux r\'eels cons\'ecutifs
        (par l'absurde, deux r\'eels cons\'ecutifs $a<b$ admettraient $c$ strictement entre eux).\hfill$\square$
\end{enumerate}

\noindent\phantomsection
{\color{BO}\bfseries\small D\'efi~6~\textendash\ \textit{Construire un syst\`eme de r\`egles logiques}}
\par\smallskip

\begin{enumerate}
  \item Applications du modus ponens~:
        \begin{itemize}
          \item \og Il pleut\fg\ + \og S'il pleut, le sol est mouill\'e\fg\ $\Rightarrow$
                \og Le sol est mouill\'e.\fg
          \item \og $n$ est divisible par $6$\fg\ + \og Si $n$ est divisible par $6$, il est
                divisible par $2$\fg\ $\Rightarrow$ \og $n$ est divisible par $2$.\fg
        \end{itemize}
  \item On a (1)~$P$, (2)~$P\Rightarrow Q$, (3)~$Q\Rightarrow R$.
        Modus ponens sur (1) et (2)~: $Q$. Modus ponens sur $Q$ et (3)~: $R$.\hfill$\square$
  \item Soit $x$ un \'el\'ement quelconque. Les hypoth\`eses \og tous les $A$ sont des $B$\fg\ et
        \og tous les $B$ sont des $C$\fg\ donnent $A(x)\Rightarrow B(x)$ et $B(x)\Rightarrow C(x)$.
        Si $A(x)$ est vrai, le modus ponens donne $B(x)$, puis $C(x)$.
        Comme $x$ est quelconque, tous les $A$ sont des $C$.\hfill$\square$
  \item Exemple~: $ABCD$ est un carr\'e $\Rightarrow$ $ABCD$ est un rectangle
        $\Rightarrow$ $ABCD$ est un parall\'elogramme $\Rightarrow$ les diagonales de $ABCD$
        se coupent en leur milieu. Si $ABCD$ est un carr\'e, on conclut par trois modus ponens.
  \item[] \textit{Extension.} Le \textbf{modus tollens} (de $P\Rightarrow Q$ et $\lnot Q$, conclure $\lnot P$)
        est la contrapos\'ee~: \og S'il a plu, le sol est mouill\'e.\fg\ \og Le sol est sec.\fg\
        $\Rightarrow$ \og Il n'a pas plu.\fg\ Le \textbf{raisonnement par cas}
        (de $(P\Rightarrow R)\land(Q\Rightarrow R)\land(P\lor Q)$, conclure $R$) est aussi valide~:
        si $P\lor Q$ est vraie, l'un des deux cas s'applique, et dans les deux $R$ est vrai.
\end{enumerate}



\needspace{3\baselineskip}
